\chapter{Multi-Agent Simulation for Insider Threat Modeling}

\section{Introduction}

Insider threats refer to security incidents that originate from within an organization and have become a critical concern in modern enterprise environments. Recent studies report that over 50\% of organizations have experienced insider incidents, with 29\% incurring remediation costs exceeding one million dollars~\cite{ponemon2025insider,gurucul2024insider}. Detecting such threats remains particularly challenging as attackers are typically trusted insiders who possess legitimate access to organizational systems. Their malicious activities often blend into normal behavior and manifest in diverse forms, including horizontal propagation such as lateral movement within internal networks, unauthorized data exfiltration, IT sabotage, fraud, and espionage, as well as vertical escalation such as privilege escalation. Well-known incidents illustrate the severity of these threats, such as Edward Snowden's data exfiltration at the NSA~\cite{harding2014snowden} and employee-driven IT sabotage at Tesla in 2018~\cite{agnihotri2024tesla}. Traditional security defenses struggle to identify insider threats due to the legitimate access privileges of insiders, and this challenge is further exacerbated by inadequate monitoring infrastructure. Recent reports show that only 36\% of organizations have deployed comprehensive monitoring systems for proactive ITD~\cite{gurucul2024insider,signpost2025insider}, revealing significant gaps in current enterprise defenses.